\chapter{Introdução}
\label{cap:introducao}

\begin{definicao}
  Um vetor do espaço é escrito na base canônica como
  $\vec{v} = a\,\ihat + b\,\jhat + c\,\khat$.
\end{definicao}

\begin{teorema}[Norma euclidiana]
  \label{teo:norma}
  Para todo $\vec{v} = a\,\ihat + b\,\jhat + c\,\khat$ vale
  \[
    \lVert \vec{v} \rVert = \sqrt{a^2 + b^2 + c^2}.
  \]
\end{teorema}

\begin{proof}
  Aplicar duas vezes o teorema de Pitágoras.
\end{proof}

\begin{exemplo}
  Se $\vec{v} = 3\,\ihat + 4\,\jhat$, então $\lVert \vec{v} \rVert = 5$
  (teorema~\ref{teo:norma}).
\end{exemplo}

\begin{observacao}
  Adicione mais capítulos na pasta \texttt{capitulos/} e inclua-os com \verb|\input|.
\end{observacao}
